\section{Related Work}

\paragraph{Context-aware speech translation.}
Document- and discourse-level context has been shown to benefit speech translation, improving quality, homophone translation, pronouns, robustness to segmentation errors, and simultaneous re-translation behavior \citep{zhang2021beyond}.  Our work differs by focusing on deployment context that is external to the current speech stream, especially slides and background documents.

\paragraph{Visual context for simultaneous translation.}
Visual context can compensate for missing source context in simultaneous translation.  \citet{caglayan2020simultaneous} show that explicit visual object-region features can improve low-latency simultaneous machine translation.  However, their setting assumes text input and image-style visual context.  We study streaming speech translation in lecture settings, where slides are temporally aligned but often asynchronous, text-heavy, and domain-specific.

\paragraph{Slides for ASR and lecture understanding.}
Slide-rich corpora and ASR studies motivate our setting.  SlideSpeech provides a large English slide-enriched audio-visual corpus with synchronized slides and transcribed speech \citep{wang2023slidespeech}.  Chinese-LiPS provides Chinese audio, lip-reading videos, slide videos, and manual transcriptions \citep{zhao2025chineselips}.  M3AV contains academic lectures with speech transcriptions, slide text, and associated papers \citep{chen2024m3av}.  Recent work asks whether slides help conference-talk transcription and finds large reductions in WER for domain-specific words when slide context is integrated \citep{sinhamahapatra2025slides}.  These works primarily target ASR or multimodal understanding rather than streaming translation with irreversible output commitments.

\paragraph{Retrieval-augmented SST and evidence timing.}
Retrieval-augmented simultaneous speech translation has recently been studied for terminology.  RASST highlights that retrieval for SST is difficult because the retriever must work under partial speech and the generator must decide whether and when to apply retrieved terms \citep{luo2026rasst}.  We generalize this idea from glossary terms to heterogeneous evidence sources including slides, background documents, and lecture history.

\paragraph{Unbounded and realistic SST.}
Recent SST systems increasingly consider long-form or unbounded streams and computation-aware latency.  InfiniSST constructs translation trajectories for unbounded speech and introduces KV-cache management to reduce computation-aware latency \citep{ouyang2025infinisst}.  Our setting is complementary: we add temporally varying external context and measure the interaction between context use, hard disambiguation, and latency.
